\begin{solapartat}
\punts{0,25} La llei de Snell dona una relació entre l'angle d'incidència i l'angle de refracció, $n_1\sin\theta_1 = n_2\sin\theta_2$, és a dir, tenim angles més grans en el medi de menor índex.

\punts{0,25} La refracció la podem tenir per angles refractats fins a $90^\circ$ ($\pi/2$~rad), ja que un angle més gran no correspon a un raig que es transmet al segon medi. Aquest angle refractat correspon a un cert angle d'incidència que anomenem angle límit. Per trobar-lo fixem $\theta_2 = \dfrac{\pi}{2}\ \mathrm{rad}$ i tenim $\sin\theta_{1c} = \dfrac{n_2}{n_1}\sin\dfrac{\pi}{2} = \dfrac{n_2}{n_1}$. A partir d'aquest angle límit sols tenim reflexió, és a dir tota la llum és reflectida, tenim reflexió total.

\figura[0.3]{sol_fig1.png}

\punts{0,25} Pel cas de la interfície aigua-aire tenim un índex del primer medi, aigua, $n_1 = 1,33$ i l'índex del segon medi és $n_2 = 1$, per tant l'angle límit $\theta_{1c}$ haurà de complir $\sin\theta_{1c} = \dfrac{n_2}{n_1} = \dfrac{1}{1,33} = 0,752$. Per tant l'angle límit: $\theta_{1c} = \arcsin(0,752) = 0,848\ \mathrm{rad} = 48,7^\circ$.

\punts{0,25} En la situació contrària tindríem els índex canviats, $n_1 = 1$ i $n_2 = 1,33$ per tant l'angle límit $\theta_{1c}$ corresponent a un angle refractat de $90^\circ$ hauria de ser més gran de $90^\circ$ cosa que no és possible.

\punts{0,25} Matemàticament ho podem comprovar en l'expressió de l'angle límit $\sin\theta_{1c} = \dfrac{n_2}{n_1}$ on es veu que necessàriament s'ha de complir $n_2 < n_1$ ja que el sinus d'un angle sempre és menor que 1.
\end{solapartat}

\begin{solapartat}
\punts{0,50} El feix corresponent a l'angle límit és la hipotenusa del triangle rectangle on els catets són l'alçada i el radi del cercle de llum.

\punts{0,25} Per tant la tangent de l'angle compleix: $\tan\theta_{1c} = \dfrac{r}{h}$.

\punts{0,25} Per tant el radi serà $r = h\tan\theta_{1c} = 2\tan 48,7^\circ = 2,28\ \mathrm{m}$

\punts{0,25} Esquema:
\figura[0.4]{sol_fig2.png}
\end{solapartat}
